\documentclass[12pt,a4paper]{article}
\usepackage{amsmath,amssymb}
\usepackage{ngerman}

\title{Laufzeit}
\author{Alexander Engelhardt, Timothy Drexler}
\date{\today}

\begin{document}
    \maketitle
    \section*{Aufgabe 3}
    \paragraph{Es gibt keinen Anlass, zu denken, die Ausführung eines Algorithmus
     auf einem Grafikprozessor beeinflusse die Komplexität jenes Algorithmus, selbst
     wenn er durch die Komputation mittels Grafikprozessors laufzeittechnisch minimal
     effizienter sei. Denn die
        ``Komplexität eines Problems ist die Laufzeit des besten Algorithmus zur Lösung dieses Problems'',
        und diese ist unabhängig vom Prozessor.}
    
    \section*{Aufgabe 4}
        \begin{enumerate}
            \item[a)] Der \textbf{Schleifenkopf} geht linear gegen n, somit ist dessen Kompl. O(n), und 
            A gegen O(1), da O(n) gr"o"ser ist, ist die Gesamtkomplexität O(n).
            \item[b)] Selbiges s.o., nur dass hier eine zweifach verschachtelte Schleife ist. Und da die
            Definition der Komplexität von k Schleifen $\mathcal{O}(n^k)$, ist hier $\mathcal{O}(n^2)$
            \item[c)] Äußere Schleife: $\mathcal{O}(n)$, innere Schleife: $\mathcal{O}(\log n)$,
            also insgesamt $n \cdot \log n \Rightarrow \mathcal{O}(n\log n)$.
            \item[d)] Iteriert gleich viele Schritte wie unter c), somit ist die Komplexität $\mathcal{O}(n\log n)$.
            \item[e)] Äußere Schleife: $\mathcal{O}(n)$, innere Schleife: $i \cdot i = n \cdot n \rightarrow \mathcal{O}(n^2)
            \Rightarrow n \cdot n^2 \\ \Rightarrow \mathcal{O}(n^3)$
            \item[f)] \begin{enumerate}
                \item Zeile 2: $\mathcal{O}(1)$.
                \item Zeile 3: ebd.
                \item Zeile 5: Teile und Herrsche mit rekursivem Funktionsaufruf: $\log n \cdot \log n = (\log n)^2 \Rightarrow \mathcal{O}(\log n)$.
                \item Zeile 6: Dadurch ebenfalls $\mathcal{O}(\log n)$.
            \end{enumerate}
            \textbf{Somit ist die maximale Komplexität der Funktion} $\mathit{O}(\log n)$.
        \end{enumerate}
\end{document}